\documentclass[17pt,aspectratio=169]{beamer}
\geometry{paperwidth=320mm,paperheight=180mm}
\setlength{\textheight}{160mm}
\usepackage{fontspec}
\setsansfont{Arial}
\setmainfont{Arial}
\usepackage{tikz,booktabs}
\definecolor{ink}{HTML}{153E38}
\definecolor{muted}{HTML}{4F625F}
\definecolor{paper}{HTML}{F8FAF6}
\definecolor{green}{HTML}{287B63}
\definecolor{pale}{HTML}{E6EFE8}
\definecolor{rust}{HTML}{A54F2A}
\setbeamercolor{normal text}{fg=ink,bg=paper}
\setbeamercolor{structure}{fg=green}
\setbeamercolor{frametitle}{fg=ink}
\setbeamercolor{block title}{fg=ink,bg=pale}
\setbeamercolor{block body}{fg=ink,bg=pale}
\setbeamerfont{frametitle}{size=\Large,series=\bfseries}
\setbeamerfont{block title}{size=\normalsize,series=\bfseries}
\setbeamersize{text margin left=9mm,text margin right=9mm}
\setbeamertemplate{navigation symbols}{}
\setbeamertemplate{itemize item}{\color{green}\rule{1.5mm}{1.5mm}}
\setbeamertemplate{footline}{\hspace{9mm}\scriptsize\color{muted}TLCN · TRÂM ANH — VƯƠNG THU — TRỌNG TÍN\hfill\insertframenumber\,/\,9\hspace{9mm}\vspace{4mm}}
\newcommand{\eyebrow}[1]{{\small\color{green}\bfseries #1}\par\vspace{2mm}}
\newcommand{\takeaway}[1]{\vfill{\color{green}\rule{\linewidth}{.5pt}}\par\vspace{1mm}{\small #1}}
\begin{document}
\begin{frame}
\eyebrow{TIỂU LUẬN CHUYÊN NGÀNH · AN TOÀN THÔNG TIN}
{\LARGE\bfseries Đã học tấn công cũ.\par\vspace{2mm}\color{green}Còn tấn công chưa từng thấy?\par}
\vspace{5mm}
{\large So sánh học có giám sát và phát hiện bất thường\par cho NIDS trên UNSW-NB15\par}
\vspace{3mm}
{\normalsize Đánh giá theo leave-one-attack-class-out\par}
\vfill
{\small Trịnh Trâm Anh · Hoàng Văn Vương Thu · Nguyễn Trọng Tín\par GVHD: ThS. Nguyễn Đăng Quang}
\end{frame}

\begin{frame}{Dữ liệu cho biết điều gì?}
\eyebrow{01 / DỮ LIỆU}
\begin{columns}[T,onlytextwidth]
\begin{column}{.47\textwidth}
{\Huge\bfseries 175.341}\par{\small dòng tập phát triển}\par\vspace{4mm}
{\Huge\bfseries 82.332}\par{\small dòng test cố định}\par
\end{column}
\begin{column}{.49\textwidth}
\begin{block}{42 feature vào mô hình}
Từ 45 cột CSV gốc, bỏ \texttt{id}, \texttt{label}, \texttt{attack\_cat}.
\end{block}
{\small Test: 37.000 Normal\par và 45.332 Attack.\par\vspace{2mm}Một dòng mô tả một flow mạng. Dữ liệu lịch sử sinh trong lab.}
\end{column}
\end{columns}
\takeaway{UNSW-NB15 là bộ dữ liệu công bố; nhóm chưa có bộ dữ liệu mạng độc lập tự thu thập.}
\end{frame}

\begin{frame}{Giấu một lớp. Giữ nguyên phép thử.}
\eyebrow{02 / THỰC NGHIỆM}
{\small CS: không giấu lớp · LR: giấu Reconnaissance\par LD: giấu DoS · LE: giấu Exploits}\par\vspace{3mm}
\begin{columns}[T,onlytextwidth]
\begin{column}{.31\textwidth}\begin{block}{1. Học}
LOACO bỏ lớp bị giấu trước khi chia train/validation.
\end{block}\end{column}
\begin{column}{.31\textwidth}\begin{block}{2. Chọn ngưỡng}
Từ Normal validation, target FPR 1\% và 5\%.
\end{block}\end{column}
\begin{column}{.31\textwidth}\begin{block}{3. Kiểm tra}
Test vẫn có lớp bị giấu; đo khả năng phát hiện.
\end{block}\end{column}
\end{columns}
\vspace{3mm}
{\Large\bfseries 4 kịch bản × 4 mô hình × 5 seed = 80 run}
\takeaway{RF/XGB học có giám sát; IF/AE học trên Normal. Test không dùng để fit hoặc chọn ngưỡng.}
\end{frame}

\begin{frame}{Một chỉ số chưa kể hết câu chuyện.}
\eyebrow{03 / KẾT QUẢ · LR–RF}
\begin{columns}[T,onlytextwidth]
\begin{column}{.48\textwidth}
{\small AP toàn test}\par
{\Huge\bfseries 0,9852}\par\vspace{5mm}
{\small AP trên Normal + Reconnaissance}\par
{\Huge\bfseries\color{rust}0,7077}
\end{column}
\begin{column}{.48\textwidth}
\begin{block}{Cùng mô hình, khác phạm vi}
Toàn test có 82.332 dòng.\par\vspace{2mm}
Normal + lớp bị giấu có 40.496 dòng.
\end{block}
{\small Trung bình 5 seed.\par AP đo khả năng xếp hạng;\par không phải tỷ lệ dự đoán đúng.}
\end{column}
\end{columns}
\takeaway{Đọc AP cùng UADR, FPR, ngưỡng và số mẫu; khác phạm vi/prevalence làm thay đổi cách hiểu.}
\end{frame}

\begin{frame}{Muốn bắt nhiều hơn, phải chịu gì?}
\eyebrow{03 / KẾT QUẢ · RF CLOSED-SET}
\begin{columns}[T,onlytextwidth]
\begin{column}{.46\textwidth}
\begin{block}{Target validation 1\%}
Recall Attack: \textbf{90,27\%}\par
FPR trên test: \textbf{4,21\%}
\end{block}
\end{column}
\begin{column}{.46\textwidth}
\begin{block}{Target validation 5\%}
Recall Attack: \textbf{97,14\%}\par
FPR trên test: \textbf{18,57\%}
\end{block}
\end{column}
\end{columns}
\vspace{4mm}
{\large Bắt được nhiều Attack hơn,\par nhưng cảnh báo nhầm Normal cũng tăng.}
\takeaway{Trung bình 5 seed. Target trên validation không bảo đảm FPR bằng target trên test hay dữ liệu mới.}
\end{frame}

\begin{frame}{Demo: đọc một cảnh báo cho đúng}
\eyebrow{04 / SUY LUẬN CSV · CS–RF–S42 · BẢN GHI 1}
{\large Nhãn thật: \textbf{Normal} \hfill Score: \textbf{0,705741}}\par\vspace{4mm}
\begin{columns}[T,onlytextwidth]
\begin{column}{.46\textwidth}
\begin{block}{Target validation 5\%}
Ngưỡng: 0,557282\par
Score $\geq$ ngưỡng\par
\textbf{\color{rust}Attack → FP}
\end{block}
\end{column}
\begin{column}{.46\textwidth}
\begin{block}{Target validation 1\%}
Ngưỡng: 0,805367\par
Score $<$ ngưỡng\par
\textbf{Normal → TN}
\end{block}
\end{column}
\end{columns}
\takeaway{Chỉ đổi ngưỡng, không học lại. FP là cảnh báo nhầm; TN là nhận đúng Normal. Không có nhãn thì chưa kết luận đúng/sai.}
\end{frame}

\begin{frame}{Chuyển sang demo khi kết nối sẵn sàng}
\eyebrow{04 / QUY TRÌNH THAO TÁC}
\begin{enumerate}
\item Chọn CS · RF · target 5\% · mẫu hỗn hợp
\item Chạy suy luận và mở bản ghi 1
\item Đổi target 1\%, chạy lại và giải thích FP → TN
\end{enumerate}
\vspace{3mm}
{\small Local: \texttt{http://127.0.0.1:4397/presentation/}\par
Cần server Python, dữ liệu và bundle trên máy.\par\vspace{2mm}
Public: \url{https://nids-unsw-nb15-tlcn.nttin181.chatgpt.site}\par
Suy luận trong trình duyệt bằng tham số đã xuất; cần tải lần đầu.\par
PDF này giữ nội dung tĩnh; không chạy CSV và không thay live demo.}
\takeaway{Khi mất mạng, dùng local đã mở sẵn hoặc ví dụ PDF. Kết nối HDMI vẫn cần thử trực tiếp tại phòng báo cáo.}
\end{frame}

\begin{frame}{Ba câu giám khảo có thể hỏi}
\eyebrow{05 / HỎI ĐÁP}
\begin{block}{“Nhóm chứng minh phát hiện mọi zero-day chưa?”}
Chưa. Nhóm thử giấu 3 lớp trên một official split.
\end{block}
\begin{block}{“CSV mới có phải dataset riêng không?”}
Chưa đủ. Cần nguồn độc lập, cách trích feature và nhãn.
\end{block}
\begin{block}{“Mô hình có đang tự học khi chạy demo không?”}
Không. Demo dùng mô hình và ngưỡng đã lưu, không fit lại.
\end{block}
\end{frame}

\begin{frame}{Đã làm được gì, còn thiếu gì?}
\eyebrow{06 / TỰ ĐÁNH GIÁ VÀ KẾT LUẬN}
\begin{columns}[T,onlytextwidth]
\begin{column}{.47\textwidth}
\begin{block}{Phần đã có}
80 run, truy vết kết quả.\par
Phân tích AP, FPR, lỗi.\par
Demo suy luận CSV.
\end{block}
{\small Giá trị: biết cách kiểm một cảnh báo và giới hạn của phép thử.}
\end{column}
\begin{column}{.49\textwidth}
\begin{block}{Phần còn mở}
Một bộ dữ liệu, còn overlap.\par
Chưa thử lưu lượng thực tế.\par
Chưa tự học hoặc ngăn chặn.
\end{block}
{\small Tiếp theo: dữ liệu ngoài, thử loại overlap, đo chi phí vận hành.}
\end{column}
\end{columns}
\takeaway{Kết luận: kết quả tốt trên benchmark cần đọc cùng phạm vi và chi phí cảnh báo nhầm trước khi triển khai thực tế.}
\end{frame}
\end{document}
